\section{Del caudal raw al caudal de bits y de entropía del puerto}
\label{sec:port-rates}
Sea \(Y\) el número de muestras raw digitalizadas por segundo, \(w\) el
ancho de cada muestra en bits y \(h\in[0,w]\) una cota inferior de
minentropía condicional por muestra válida para el modelo indicado abajo.
Los dos caudales brutos son \(B_0=wY\) bit/s y \(X=hY\) bits de
minentropía/s. Si se conocen directamente \(X,B_0,w\), se introduce
\(Y=B_0/w\) y \(h=wX/B_0\) para \(B_0>0\).
Sólo con los dos caudales globales es imposible determinar un único resultado.
Por ejemplo, si los datos destinados al selector concentran la aleatoriedad
y las ventanas de datos son constantes, la fuente puede tener entropía
positiva mientras la salida Toeplitz es constante.

\paragraph{Modelo del reparto de muestras.}
Por intento se asignan \(n\) muestras de datos, \(u\) muestras a renovar el
selector y \(g\) muestras de guarda. La matriz es de
\(m\times wn\), con \(m=128\), y su índice tiene \(d=wn+m-1\) bits.
Si se consumen \(t\) incrementos de \(d\) bits empaquetados en muestras
de \(w\) bits, \(u=t\lceil d/w\rceil\). Reservar esas muestras permite
producir un índice de esa anchura; su calidad y su relación con los datos
se justifican por separado. Sea \(p\) la fracción asintótica de intentos
admitidos por las pruebas de la fuente, en un régimen estable con caudal
\(Y\). Si todos los intentos cuestan \(C=n+u+g\), el caudal de números
intermedios aceptados es \(\nu=pY/C\) bloques/s antes de aplicar límites de
cómputo. Con paradas o rearranques se incluye su tiempo en el denominador.
\[
\begin{array}{lcc}
\text{Etapa, a plena adquisición} & \text{bits producidos/s}
   & \text{minentropía/s destinada a datos}\\
\hline
\text{Digitalizador} & wY & X=hY\ \text{ (bruta)}\\
\text{Ventanas de datos aceptadas} & wn\nu & k\nu,\quad k=nh\\
\text{Índices de lotes aceptados} & d\nu & 0\ \text{ imputados al puerto}\\
\text{PTG.2 tras Toeplitz} & m\nu & \kappa\nu
\end{array}
\]
Las filas representan etapas consecutivas. Para aplicar \(k=nh\) se requiere
una cota de bloque bajo las condiciones de aceptación, historial y vista
declaradas. Puede demostrarse directamente para la ventana sin suponer
muestras IID. Si sólo se conoce la entropía de las muestras individuales,
el producto \(nh\) requiere una justificación adicional. Los índices, los
rechazos y las guardas no aportan crédito de entropía de salida.
Para conciliar el contador de entrada con los lotes aceptados, los intentos
rechazados consumen \((1-p)wY\) bit/s. En los aceptados se destinan
\(wu\nu\) bit/s a producir los índices y \(wg\nu\) bit/s a guardas.
En particular, \(d\nu\) cuenta los bits de índice disponibles después de
renovarlo; las muestras consumidas para esa renovación pueden ocupar más.

\paragraph{Qué significa la columna de entropía.}
Para dar una tasa \emph{sin suavizado} en el caso clásico, supóngase que,
para todo historial admisible y valor de la vista concedida, la masa máxima
de la ventana raw aceptada en el modelo ideal es a lo sumo \(2^{-k}\) y
que la distribución conjunta real del índice y la ventana está a distancia
a lo sumo \(\eta\) de un índice uniforme independiente y esa ley ideal.
La vista para estas hipótesis contiene lo suficiente para fijar el estado
anterior, el historial de entradas, los tiempos y alertas observables y,
en la ruta con CMAC, los prefijos usados. El índice actual permanece interno.
Estas cotas uniformes en el contexto son más fuertes que una cota promedio
sobre las realizaciones del entorno.

Con \(P_u\) la distribución ideal de la salida a índice fijado, la
universalidad de Toeplitz da
\(\mathbb E_u\lVert P_u-U_m\rVert_2^2\le(1-2^{-m})2^{-k}\).
Para cada índice, la mayor masa es a lo sumo
\(2^{-m}+\lVert P_u-U_m\rVert_2\). Al promediar y aplicar
Cauchy--Schwarz, y trasladar cualquier suceso de adivinación a la ley real
con coste \(\eta\), se obtiene la cota por bloque
\[
 \kappa=-\log_2\min\!\left\{1,\,
 2^{-m}+\sqrt{(1-2^{-m})2^{-k}}+\eta\right\}.
\]
Ocultar el índice permite usar este promedio por cada historial fijado.
Si se concede la secuencia de índices, se han de rejustificar los
condicionamientos frente a esa nueva vista. Como la cota por bloque vale
tras cualquier prefijo admisible y para cada valor de la información
lateral clásica concedida, la regla del producto acota cada secuencia de
\(q\) salidas por \(2^{-q\kappa}\). Las biyecciones de ambos núcleos
transfieren la garantía a la misma vista. Si el modelo físico respalda
directamente otra cota puntual \(\kappa_{\rm fuente}\), se puede insertar
esa cota en las fórmulas de tasa. El valor AIS de 0.98 por bit de número
intermedio no acredita por sí solo la premisa conjunta
\cite[§§335,382 y 402]{ais}.

Para una garantía ITS en distancia se conserva por separado
\(\epsilon_q\le\min\{1,q(2^{(m-k)/2-1}+2\eta)\}\) bajo las
hipótesis del LHL y la composición del paper. La distancia a uniforme
permite \(qm\) bits de minentropía \emph{suavizada} en distancia de
variación de radio \(\epsilon_q\) para la ley clásica. Si sólo se dispone
de una distancia conjunta \(\epsilon_q\), para la minentropía
\emph{sin suavizado} se obtiene
\(H_{\min}(Y^q\mid E)\ge -\log_2\min(1,2^{-mq}+\epsilon_q)\).
Así se evita multiplicar un certificado marginal por el número de bloques.
Para información lateral cuántica, el teorema de permutación sigue siendo
aplicable: la tasa estricta utiliza una cota \emph{conjunta} en el modelo
cq. La derivación puntual clásica anterior tiene su propio alcance.

\subsection{Puerto AES/SHA: reserva de siembra y conservación 1:1}
\label{sec:port-rate-ais}
Sea \(L\) el número de peticiones de 128 bits por época. El puerto asigna
dos bloques del PTG.2 a la siembra o resiembra y \(L\) bloques a datos.
Con \(f_A=L\nu/(L+2)\) peticiones/s, el reparto de los \(m\nu\)
bits/s intermedios es \(2m\nu/(L+2)\) hacia la semilla y \(mf_A\) al
puerto de datos. La permutación en la entrada fresca entrega
\[
 B_A=128 f_A,\qquad E_A\ \ge\ \kappa f_A.
\]
Aquí y en CTR, \(E_A\) o \(E_C\) es el crédito de minentropía conjunta
por tiempo transcurrido a caudal estable. En una traza parada se contabilizan
sus bloques realmente entregados y su duración. Si el número de entregas o
la causa de parada se divulgan y dependen de los datos, la premisa puntual
debe mantenerse también condicionada a esa información. Sustituir un número
aleatorio de entregas por su esperanza en una cota de minentropía conjunta
requeriría otro argumento. El consumo del estado SHA
no se descuenta como una segunda adquisición. El presupuesto de una sesión
finita suma el calentamiento del selector y los reinicios efectivos.

\subsection{Puerto CTR: coste de los tres puertos de material físico}
\label{sec:port-rate-ctr}
Contemos \(s\) bloques Toeplitz para la siembra o resiembra por época y
\(c\) bloques adquiridos por palabra devuelta. El puerto CTR directo inyecta
\(0^{256}\|A\) y entrega \(m\) bits aunque la representación de la
entrada adicional ocupe 384. El caudal de peticiones al
puerto es \(f_C=L\nu/(s+cL)\); las ramas de adquisición de
semilla y de datos reciben, respectivamente, \(ms\nu/(s+cL)\) y
\(mcL\nu/(s+cL)\) bits/s. En el perfil directo con CMAC por petición,
dos bloques aportan los dos operandos del CMAC y se entrega un tag al CTR.
\[
 B_C=m f_C,\qquad E_C\ \ge\ \kappa f_C .
\]
Para mover esta misma cota a través de CMAC se exige \(\kappa\) en el
último operando condicionado al primero y al resto de la vista.
La validación normativa requiere además el crédito por llamada indicado en
la sección~\ref{sec:puertos-nist}. En RS la coordenada conservada es el
último bloque del material de resiembra, condicionado a los 256 primeros
bits y al estado. En XOR la coordenada física debe satisfacer la misma
premisa condicionada a la máscara del DRBG. Las asignaciones para una
petición externa de 128 bits son
\[
\begin{array}{lrr}
\text{Diseño y asignación de material} & s & c\\
\hline
\text{CTR directo, puerto de siembra validado full entropy} &3&1\\
\text{CTR directo, CMAC sólo para siembra} &6&1\\
\text{CTR directo, CMAC también para cada salida} &6&2\\
\text{RBG3(XOR), CMAC en siembra y en rama física} &6&2\\
\text{RBG3(RS), CMAC en siembra y en cada reseed} &6&6
\end{array}
\]
Para RS, los \(s\) bloques iniciales se cargan una vez al instanciar. Cada
iteración RS genera 256 bits internos, \(256 f_C\) bit/s bajo esta
contabilidad, y trunca a 128. Los números de coste de la tabla presuponen
los componentes y autorizaciones de la sección NIST; describen varios
puertos con reclamaciones distintas. Cada tag completo de 128 bits en la
ruta externa vetted requiere al menos 192 bits de entropía acreditada por
evaluación. La reserva \(s=3\) en el perfil directo requiere un puerto
autorizado de entropía completa. La identidad ITS por sí sola no lo valida.

\paragraph{Tasa a la salida del equipo.}
Sean \(T_{\rm acq}=(s+cL)/\nu\), \(T_T=(s+cL)/F_T\), con \(F_T\)
en bloques Toeplitz/s, y \(T_D=128L/B_D+T_{\rm aux}\), con \(B_D\)
en bits/s de núcleo y \(T_{\rm aux}\) en segundos para el trabajo auxiliar
de la época sobre ese mismo recurso. AES/SHA usa \((s,c)=(2,1)\).
En ejecución serial, \(T=T_{\rm acq}+T_T+T_D\). Para etapas independientes
solapadas con capacidad de búfer suficiente, el tiempo por época en régimen
estacionario ideal es \(T=\max(T_{\rm acq},T_T,T_D)\). Si comparten un
recurso se suman sus cargas antes de tomar el máximo. Tiempos de llenado o
vaciado corresponden al cálculo de una ráfaga finita. Los resultados finales
son \(B_{\rm puerto}=128L/T\) y \(E_{\rm puerto}\ge\kappa L/T\).
El estado de fallo persistente lleva ambos caudales a cero.

\paragraph{Ejemplo común reproducible.}
Adóptense \(Y=10^8\), \(w=1\), \(h=0.75\), \(n=427\), \(t=2\),
\(g=32\), \(p=0.99\), \(L=1024\) y, exclusivamente como premisa sintética por
historial, \(\eta=2^{-160}\). Resultan \(B_0=100\) Mbit/s,
\(X=75\) Mbit de minentropía/s, \(k=320.25\), \(d=554\),
\(C=1567\), \(\nu=63178.05\) bloques/s y
\(\kappa\ge127.99999999935\) bits por bloque bajo las premisas anteriores.
Los datos aceptados representan 26.9770 Mbit/s y al menos 20.2328 Mbit
de minentropía/s antes de extraer; los índices aceptados suman 35.0006
Mbit/s de representación interna. El PTG.2 entrega 8.08679 Mbit/s.
Para \(q=2^{32}\) bloques de datos, la cota de distancia es inferior a
\(2^{-64}\). Si la afirmación de distancia incluye además los bloques
de semilla se usa el número total de extracciones pertinente.

Con capacidad Toeplitz suficiente, sin coste auxiliar añadido y las tasas
de núcleo software medidas, los puertos directos AES/SHA y CTR alcanzan en
la tubería modelada 8.071026 y 8.063167 Mbit/s respectivamente. En este
ejemplo ambas cotas estrictas de entropía/s coinciden con los caudales de
bits/s a los seis decimales mostrados; mantienen el factor exacto
\(\kappa/128\). En operación serial con tiempo Toeplitz idealizado a cero
se obtienen 6.14935 y 4.23232 Mbit/s. Para CTR con CMAC sólo en la
semilla, el límite de adquisición baja a 8.03968 Mbit/s. Si se usa además
CMAC en cada salida baja a 4.03158 Mbit/s, y en RS con CMAC a 1.34648
Mbit/s para sesiones de 1024 entregas. Romper las premisas de entropía o
dependencia puede invalidar la columna de minentropía sin cambiar el
contador de bits. Los cálculos son reproducibles con
\path{experiments/rate_budget.py}. La opción \path{--raw-bit-s} recibe
el caudal raw en bits/s y \path{--raw-entropy-bit-s} la cota de entropía
raw en bits/s. El JSON muestra también la pérdida de entropía por bloque
para resolver diferencias pequeñas en la aritmética de coma flotante.
